\section*{Chapter \ref{ch:m3:freight-weather-and-other-corners} --- Freight, Weather and Other Corners}

\begin{solution}{exo:m3:freight-weather-and-other-corners:1}
The average is 31°F: 34 HDD and 0 CDD.
\end{solution}

\begin{solution}{exo:m3:freight-weather-and-other-corners:2}
$(1\,500\,000 - 450\,000 - 100\,000)/38 = \$25\,000$ a day.
\end{solution}

\begin{solution}{exo:m3:freight-weather-and-other-corners:3}
An average is harder to move by one day's assessment, matches the way ships are fixed through a month, and
makes the index hedgeable by fixing several ships over the period; it also reduces the effect of one broker
panel's judgement on a single day.
\end{solution}

\begin{solution}{exo:m3:freight-weather-and-other-corners:4}
The index is 646 below the strike: the buyer (long HDD) pays the seller $646 \times 20\,000 = \$12.92$
million.
\end{solution}

\begin{solution}{exo:m3:freight-weather-and-other-corners:5}
It receives $(21\,000 - 19\,500) \times 30 = \$45\,000$, which offsets the lower hire it earns in the market:
it has fixed an average of \$21\,000 a day for the month, up to the basis between its own fixture and the
index average.
\end{solution}

\begin{solution}{exo:m3:freight-weather-and-other-corners:6}
The trend is negative: winters are getting warmer. Moving old winters onto the 2026/27 trend makes them warmer
than they were, so more of them fall below the strike and the average payout rises.
\end{solution}

\begin{solution}{exo:m3:freight-weather-and-other-corners:7}
36 complete winters, 1990/91 to 2025/26; mean 4\,912 HDD, standard deviation 477; warmest 2011/12 (3\,932),
coldest 2013/14 (5\,999); trend $-125$ HDD a decade.
\end{solution}

\begin{solution}{exo:m3:freight-weather-and-other-corners:8}
Black--Scholes prices by replication in a traded underlying; an HDD index is not traded and cannot be held or
hedged continuously, its distribution is bounded, seasonal and trend-affected, and the season's index is
known only as it accumulates. Price from the index's own distribution (\omterm{def:m3:freight-weather-and-other-corners:burn}{burn analysis} or simulation of daily
temperatures) plus a risk margin.
\end{solution}

\begin{solution}{pb:m3:freight-weather-and-other-corners:1}
\textbf{1.} 36 winters, mean 4\,912 HDD. \textbf{2.} 477 HDD. \textbf{3.} Warmest 2011/12 with 3\,932,
coldest 2013/14 with 5\,999. \textbf{4.} $-125$ HDD a decade. \textbf{5.} Recent winters are warmer; an
undetrended average overstates the next winter's HDD and underprices the put. \textbf{6.} 4\,677 HDD.
\textbf{7.} 4\,209. \textbf{8.} 20\% (six of thirty winters). \textbf{9.} \$0.61 million on average;
\$1.60 million in the 90th-percentile winter. \textbf{10.} \$9.30 million, in the winter of 2011/12 adjusted to
the 2026/27 trend. \textbf{11.} The seller holds a skewed risk it cannot hedge: it charges for capital, for the
uncertainty in the trend and for the thin history. \textbf{12.} The seller's worst case; the utility gives up
protection beyond \$20 million. \textbf{13.} The utility's margin depends on its customers' temperatures and
behaviour, not exactly O'Hare's HDD: station and volume basis. \textbf{14.} A swap costs nothing upfront but gives
up the gain of a cold winter; the put keeps it for a premium. Shareholders may prefer stable earnings (swap) or
upside (put). \textbf{15.} Energy traders with opposite exposures, reinsurers, and funds seeking uncorrelated risk;
they diversify across stations and seasons. \textbf{16.} Tail winters rarer than one in thirty, changes at the
station, and the trend's own uncertainty. \textbf{17.} A forecast of a mild winter raises the put's price, as the
distribution shifts; closer to the season the forecast outweighs history. \textbf{18.} Hedgers' exposures are
local and hard to match with a few stations; basis is large and there are few natural sellers.
\textbf{19.} \emph{Named result:} a detrended burn estimate of 4\,677 HDD, a strike of 4\,209, and a
one-in-ten-year (90th percentile) payout of \$1.60 million. \textbf{20.} A contract that pays on a measured weather
index at a named place, so that those whose income depends on the weather can transfer the risk.
\end{solution}

\begin{solution}{iq:m3:freight-weather-and-other-corners:1}
$\max(0, 65 - \theta)$ with $\theta$ the day's mean temperature in °F, summed over the season. Buyers of
protection against warm winters: gas and power utilities, heating-oil distributors; against cold: those whose
costs rise with heating demand.
\iqlookfor{the definition and the natural hedgers.}
\end{solution}

\begin{solution}{iq:m3:freight-weather-and-other-corners:2}
A cash-settled forward on a freight \omterm{def:m3:freight-weather-and-other-corners:route}{route index} (dollars per day or per tonne) for a future period, settling
on the average of the index over that period against the agreed rate.
\iqlookfor{cash settlement on an average of a panel index.}
\end{solution}

\begin{solution}{iq:m3:freight-weather-and-other-corners:3}
From the index's distribution: \omterm{def:m3:freight-weather-and-other-corners:burn}{burn analysis} on detrended history, or simulation of daily temperatures with
seasonality and autocorrelation, plus a risk loading. Not Black--Scholes: there is no traded underlying to
replicate with, and the index is bounded and seasonal.
\iqlookfor{no replication; actuarial plus risk premium.}
\end{solution}

\begin{solution}{iq:m3:freight-weather-and-other-corners:4}
Sell FFAs on the relevant route for the months after delivery (or sell a \omterm{def:m3:freight-weather-and-other-corners:charter}{time charter} now). Remaining risk:
basis between its ship and the index route, timing, and the ship's availability.
\iqlookfor{the hedge and its basis.}
\end{solution}

\begin{solution}{iq:m3:freight-weather-and-other-corners:5}
A short position in warm-winter tails, correlated across cities through large-scale patterns. Measure it with
joint historical and simulated winters across the stations, capped payouts, and stress scenarios such as an
El Niño winter.
\iqlookfor{correlation across stations and tail measures.}
\end{solution}

\begin{solution}{iq:m3:freight-weather-and-other-corners:6}
Ingest station data with quality flags; fill nothing silently; compute indices by the contract's rules;
detrend and simulate; blend forecasts for near seasons; price payoffs; publish quotes with the risk loading;
monitor station changes and data revisions.
\iqlookfor{data quality, contract rules, model and forecast blending.}
\end{solution}
